% SGA3.c.pdf，PDF 第63--86页；止于第5节之前。
\SourcePage{63}
\SGAchapter{XXI}{根数据}{Données radicielles}
\begin{center}M. Demazure\end{center}

由于缺少适当的参考文献，\footnote{编者注：关于\Term{root-system}{根系}的结果（第1--5节），可参阅 [BLie]，第 VI 章。}本讲汇集关于\Term{root-datum}{根数据}（即“抽象”\Term{root-system}{根系}）的一些已知结果，其中大部分将在后文用到。

\textbf{记号。}以 $\mathbb Q_+$ 表示正有理数（包括零）的\Term{set}{集合}；于是 $\mathbb Z\cap\mathbb Q_+=\mathbb N$。设 $V$ 是一个 $\mathbb Q$-\Term{vector-space}{向量空间}；若 $A$（分别地，$B$）是 $\mathbb Q$（分别地，$V$）的\Term{subset}{子集}，以 $A\cdot B$ 表示 $A\otimes B$ 在\Term{morphism}{态射} $\mathbb Q\otimes V\to V$ 下的像，也就是说，它是 $B$ 中元素的系数在 $A$ 中的线性组合所成的\Term{set}{集合}。记 $-B=\{-1\}B$。以 $E-F$ 表示属于 $E$ 而不属于 $F$ 的元素的\Term{set}{集合}。

\SGAsection{1}{一般性质}
\SGAsubsection{1.1}{定义与初步性质}
\begin{SGAresult}{定义}{1.1.1}
设 $M$ 与 $M^*$ 是两个互相\Term{duality}{对偶}的\Term{finite-type}{有限型}自由 $\mathbb Z$-模。记 $V=M\otimes\mathbb Q$，$V^*=M^*\otimes\mathbb Q$；它们是两个互相\Term{duality}{对偶}的 $\mathbb Q$-\Term{vector-space}{向量空间}。将 $M$（分别地，$M^*$）视为 $V$（分别地，$V^*$）的\Term{subset}{子集}。以 $(\ ,\ )$ 表示 $M^*\times M$（分别地，$V^*\times V$）上的\Term{canonical}{典范}\Term{bilinear}{双线性}\Term{form}{形式}。

设 $R$ 是 $M$ 的一个\Term{finite}{有限}\Term{subset}{子集}。给定从 $R$ 到 $M^*$ 的映射 $\alpha\mapsto\alpha^*$；以 $R^*$ 表示所有 $\alpha^*$（$\alpha\in R$）所成的\Term{set}{集合}。对每个 $\alpha\in R$，用下列公式定义 $M$ 和 $V$（分别地，$M^*$ 和 $V^*$）的\Term{endomorphism}{自同态} $s_\alpha$（分别地，$s_\alpha^*$）：
\begin{align}
s_\alpha(x)&=x-(\alpha^*,x)\alpha,\qquad\text{即 }s_\alpha=\operatorname{id}-\alpha^*\otimes\alpha;\tag{1}\\
s_\alpha^*(u)&=u-(u,\alpha)\alpha^*,\qquad\text{即 }s_\alpha^*=\operatorname{id}-\alpha\otimes\alpha^*.\tag{$1^*$}
\end{align}
若满足下列公理，就说二元组 $(R,R^*)$（更准确地说，二元组 $(R,R\to M^*)$）是 $(M,M^*)$ 中的一个\Term{root-datum}{根数据}，或说 $(M,M^*,R,R^*)$ 是一个\Term{root-datum}{根数据}：
\SourcePage{64}
\begin{align*}
\text{(DR I)}&\quad\text{对每个 }\alpha\in R,\quad (\alpha^*,\alpha)=2;\\
\text{(DR II)}&\quad\text{对每个 }\alpha\in R,\quad s_\alpha(R)\subset R,\quad s_\alpha^*(R^*)\subset R^*.
\end{align*}
称 $R$ 为\Term{root-datum}{根数据} $\mathscr R=(M,M^*,R,R^*)$ 的\Term{root-system}{根系}。$R$（分别地，$R^*$）中的元素称为该\Term{root-datum}{根数据}的\Term{root}{根}（分别地，\Term{coroot}{余根}）。
\end{SGAresult}

\begin{SGAresult}{注记}{1.1.2}
公理 (DR I) \Term{equivalence}{等价}于下列任一性质：
\begin{align*}
\text{(2)}\quad s_\alpha s_\alpha&=\operatorname{id},&\text{($2^*$)}\quad s_\alpha^*s_\alpha^*&=\operatorname{id},\\
\text{(3)}\quad s_\alpha(\alpha)&=-\alpha,&\text{($3^*$)}\quad s_\alpha^*(\alpha^*)&=-\alpha^*.
\end{align*}
\end{SGAresult}
\begin{SGAresult}{注记}{1.1.3}
公理 (DR I) 和 (DR II) 蕴涵
\[R=-R,\qquad R^*=-R^*,\qquad 0\notin R,\qquad0\notin R^*.\]
\end{SGAresult}
\begin{SGAresult}{引理}{1.1.4}
映射 $R\to R^*$ 是\Term{bijection}{双射}。更一般地，若 $\alpha,\beta\in R$ 且对每个 $x\in R$ 都有 $(\alpha^*,x)=(\beta^*,x)$，则 $\alpha=\beta$。
\end{SGAresult}
\begin{proof}
事实上，此时 $s_\beta(\alpha)=\alpha-2\beta$，$s_\alpha(\beta)=\beta-2\alpha$。立即得到
\[
s_\beta s_\alpha(\alpha)=2\beta-\alpha=\alpha+2(\beta-\alpha),\qquad
s_\beta s_\alpha(\beta-\alpha)=s_\beta(\beta-\alpha)=\beta-\alpha,
\]
因此，由 (DR II)，$(s_\beta s_\alpha)^n(\alpha)=\alpha+2n(\beta-\alpha)\in R$。由于 $R$ \Term{finite}{有限}，故 $\beta-\alpha=0$。
\end{proof}
\begin{SGAresult}{推论}{1.1.5}
逆映射 $R^*\to R$ 定义一个\Term{root-datum}{根数据}
\[\mathscr R^*=(M^*,M,R^*,R),\]
称为 $\mathscr R$ 的\Term{dual-root-datum}{对偶根数据}。\footnote{编者注：注意 $(\alpha^*)^*=\alpha$。}
\end{SGAresult}
\begin{SGAresult}{定义}{1.1.6}
以 $\Gamma_0(R)$ 表示 $M$ 中由 $R$ 生成的\Term{subgroup}{子群}。以 $\mathscr V(R)$ 表示 $V$ 中由 $R$ 张成的向量子\Term{space}{空间}，即 $\Gamma_0(R)\otimes\mathbb Q$。把这些定义应用于 $\mathscr R^*$，同样得到 $\Gamma_0(R^*)$ 与 $\mathscr V(R^*)$。

称下列数为 $\mathscr R$ 的\Term{reductive-rank}{约化秩}：
\[\operatorname{rgred}(\mathscr R)=\operatorname{rang}(M)=\dim(V)=\dim(V^*)=\operatorname{rang}(M^*)=\operatorname{rgred}(\mathscr R^*).\]
称下列数为 $\mathscr R$ 的\Term{semisimple-rank}{半单秩}：
\[\operatorname{rgss}(\mathscr R)=\operatorname{rang}(R)=\operatorname{rang}(\Gamma_0(R))=\dim(\mathscr V(R)).\]
因此 $\operatorname{rgss}(\mathscr R)\leq\operatorname{rgred}(\mathscr R)$。

下文将证明 $\operatorname{rgss}(\mathscr R)=\operatorname{rgss}(\mathscr R^*)$，也就是说，$\mathscr V(R)$ 与 $\mathscr V(R^*)$ 的\Term{dimension}{维数}相同。
\end{SGAresult}
\begin{SGAresult}{定义}{1.1.7}
若 $\operatorname{rgss}(\mathscr R)=\operatorname{rgred}(\mathscr R)$（分别地，$\operatorname{rgss}(\mathscr R)=0$），称 $\mathscr R$ 是\Term{semisimple}{半单}的（分别地，\Term{trivial-root-datum}{平凡}的）。所以 $\mathscr R$ \Term{trivial-root-datum}{平凡}的充要条件是 $R$ 为空。\Term{reductive-rank}{约化秩}为零的\Term{trivial-root-datum}{平凡根数据}记为 $\mathbf0=(\{0\},\{0\},\varnothing,\varnothing)$。
\end{SGAresult}

\SourcePage{65}
\begin{SGAresult}{定义}{1.1.8}
以 $W(\mathscr R)$ 表示 $M$ 的变换群中由 $s_\alpha$（$\alpha\in R$）生成的群，称为 $\mathscr R$ 的\Term{weyl-group}{Weyl 群}。记
\[W^*(\mathscr R)=W(\mathscr R^*).\]
于是 $W(\mathscr R)$ \Term{action}{作用}于 $R,\Gamma_0(R),\mathscr V(R),M,V$。若 $w\in W(\mathscr R)$ 且 $x\in M$（分别地，$x\in V$），则 $wx-x\in\Gamma_0(R)$（分别地，$wx-x\in\mathscr V(R)$），这由公式 (1) 立即可见。$W^*(\mathscr R)$ 也有相同性质。
\end{SGAresult}
\begin{SGAresult}{引理}{1.1.9}
对所有 $\alpha\in R$、$x\in V$、$u\in V^*$，有
\[(s_\alpha^*(u),s_\alpha(x))=(u,x).\tag{4}\]
\end{SGAresult}
\begin{proof}
事实上，对 (1) 与 ($1^*$) 作内积，得到左端等于
\[(u,x)+(u,\alpha)(\alpha^*,x)\bigl((\alpha^*,\alpha)-2\bigr)=(u,x).\]
\end{proof}
\begin{SGAresult}{注记}{1.1.10}
若假设 $0\notin R$ 且 $0\notin R^*$，则 1.1.9 \Term{equivalence}{等价}于 (DR I)。
\end{SGAresult}
\begin{SGAresult}{推论}{1.1.11}
以 $h\mapsto h^\vee$ 表示从 $\operatorname{GL}(M)$ 到 $\operatorname{GL}(M^*)$ 的\Term{isomorphism}{同构}，它把 $h$ 映到其\Term{contragredient}{逆步变换}。则公式 (4) 也可写为
\[(s_\alpha)^\vee=s_\alpha^*.\tag{5}\]
\end{SGAresult}
\begin{SGAresult}{推论}{1.1.12}
上述\Term{isomorphism}{同构}诱导 $W(\mathscr R)$ 到 $W^*(\mathscr R)$ 的一个\Term{isomorphism}{同构}。
\end{SGAresult}
\begin{SGAresult}{说明}{1.1.13}
根据上一结果，我们将 $W$ 与 $W^*$ 等同，并把 $W$ 看作 $R$、$R^*$、$M$、$M^*$、$\Gamma_0(R)$、$\Gamma_0(R^*)$、$V$、$V^*$、$\mathscr V(R)$、$\mathscr V(R^*)$ 的变换群。以后用 $s_\alpha$ 代替 $s_\alpha^*$。
\end{SGAresult}
\SGAsubsection{1.2}{映射 $p$}
\begin{SGAresult}{引理}{1.2.1}
设 $p:M\to M^*$（分别地，$V\to V^*$）为下式定义的线性映射：
\[p(x)=\sum_{u\in R^*}(u,x)u.\tag{6}\]
记 $\ell(x)=(p(x),x)$。有下列性质：
\begin{align}
\ell(x)&\geq0,\qquad\ell(\alpha)>0\quad(\alpha\in R).\tag{7}\\
\ell(wx)&=\ell(x)\quad(w\in W).\tag{8}\\
(p(x),y)&=(p(y),x)\quad(x,y\in V).\tag{9}\\
\ell(\alpha)\alpha^*&=2p(\alpha)\quad(\alpha\in R).\tag{10}
\end{align}
\end{SGAresult}
\SourcePage{66}
\begin{proof}
前三个关系显然成立。\footnote{编者注：(8) 由 (4)、(2) 和 (DR II) 得出。}下面证明最后一个。由 (1)，对 $u\in V^*$ 有
\begin{align*}
(u,\alpha)^2\alpha^*&=(u,\alpha)u-(u,\alpha)s_\alpha(u)\\
&=(u,\alpha)u+(u,s_\alpha(\alpha))s_\alpha(u)\\
&=(u,\alpha)u+(s_\alpha(u),\alpha)s_\alpha(u)\qquad(\text{因为 }s_\alpha^2=\operatorname{id}).
\end{align*}
由 (DR II)，$s_\alpha$ 是 $R^*$ 的一个置换，因此只须对 $u\in R^*$ 求和即可。
\end{proof}
\begin{SGAresult}{说明}{1.2.2}
关系 (10) 表明 $\alpha\mapsto\ell(\alpha)\alpha^*$ 是 $M$ 到 $M^*$ 的一个线性映射在 $R$ 上的\Term{restriction}{限制}。特别地，
\[(-\alpha)^*=-\alpha^*.\]
\end{SGAresult}
\begin{SGAresult}{推论}{1.2.3}
映射 $p$ 诱导 $\mathscr V(R)$ 到 $\mathscr V(R)^*$ 的一个\Term{isomorphism}{同构}。
\end{SGAresult}
\begin{proof}
事实上，$p$ 将 $\mathscr V(R)$ \Term{surjectivity}{满射}到 $\mathscr V(R^*)$。因此
\[\dim(\mathscr V(R))\geq\dim(\mathscr V(R^*)).\]
把此式应用于\Term{dual-root-datum}{对偶根数据}，得到
\[\dim(\mathscr V(R))=\dim(\mathscr V(R^*)),\]
故 $p$ 既为\Term{surjectivity}{满射}，也为\Term{bijection}{双射}。
\end{proof}
\begin{SGAresult}{推论}{1.2.4}
有 $\dim(\mathscr V(R))=\dim(\mathscr V(R^*))$，从而
\[\operatorname{rgss}(\mathscr R)=\operatorname{rgss}(\mathscr R^*).\]
\end{SGAresult}
\begin{SGAresult}{推论}{1.2.5}
\Term{bilinear}{双线性}\Term{form}{形式} $(u,x)$ 在 $\mathscr V(R^*)\times\mathscr V(R)$ 上非\Term{degeneration}{退化}，因而使这两个 $\mathbb Q$-\Term{vector-space}{向量空间}互相\Term{duality}{对偶}。
\end{SGAresult}
\begin{proof}
事实上，若对每个 $u\in R^*$ 都有 $(u,x)=0$，则 $p(x)=0$。
\end{proof}
\begin{SGAresult}{推论}{1.2.6}
对称\Term{bilinear}{双线性}\Term{form}{形式} $(p(x),y)$ 在 $\mathscr V(R)$ 上正定。
\end{SGAresult}
\begin{SGAresult}{推论}{1.2.7}
$W$ \Term{faithful-action}{忠实地作用}于 $R$（因而也忠实地\Term{action}{作用}于 1.1.13 中的其他\Term{set}{集合}）。
\end{SGAresult}
\begin{proof}
事实上，设 $u\in V^*$。设 $w\in W$，并假设对每个 $\alpha\in R$ 都有 $w(\alpha)=\alpha$，要证明 $w(u)=u$。有
\[(w(u)-u,\alpha)=(w(u),\alpha)-(u,\alpha)=(u,w^{-1}(\alpha))-(u,\alpha)=0.\]
但 $w(u)-u\in\mathscr V(R^*)$。若它与所有\Term{root}{根}都\Term{orthogonality}{正交}，则由 1.2.5，它为零。
\end{proof}
\begin{SGAresult}{推论}{1.2.8}
群 $W$ 是\Term{finite-group}{有限群}。
\end{SGAresult}

\SourcePage{67}
\begin{SGAresult}{命题}{1.2.9}
$W$ 的\Term{action}{作用}保持\Term{root}{根}与\Term{coroot}{余根}之间的对应。换言之，对 $\alpha\in R$ 与 $w\in W$，有
\[w(\alpha^*)=w(\alpha)^*.\]
\end{SGAresult}
\begin{proof}
只须在 $w=s_\beta$（$\beta\in R$）时证明，即验证公式
\[s_\beta(\alpha^*)=s_\beta(\alpha)^*.\]
而 $\ell(s_\beta(\alpha))s_\beta(\alpha)^*/2$ 等于
\begin{align*}
p(s_\beta(\alpha))&=\sum_{u\in R^*}(u,s_\beta(\alpha))u
=\sum_{u\in R^*}(s_\beta(u),\alpha)u\\
&=\sum_{u\in R^*}(u,\alpha)s_\beta(u)=s_\beta(p(\alpha));
\end{align*}
因 $\ell(s_\beta(\alpha))=\ell(\alpha)$，便得到 $s_\beta(\alpha)^*=s_\beta\bigl(2p(\alpha)/\ell(\alpha)\bigr)=s_\beta(\alpha^*)$。
\end{proof}
\begin{SGAresult}{推论}{1.2.10}
若 $\alpha\in R$ 且 $w\in W$，则
\[ws_\alpha w^{-1}=s_{w(\alpha)}.\]
\end{SGAresult}
\begin{proof}
事实上，$ws_\alpha w^{-1}(x)=x-(\alpha^*,w^{-1}(x))w(\alpha)=x-(w(\alpha^*),x)w(\alpha)$，而由 1.2.9，最后一项等于
\[x-(w(\alpha)^*,x)w(\alpha)=s_{w(\alpha)}(x).\]
\end{proof}
\begin{SGAresult}{推论}{1.2.11}
设 $R'\subset R$。$(M,M^*,R',R'^*)$ 为\Term{root-datum}{根数据}的充要条件是：$\alpha,\beta\in R'$ 蕴涵 $s_\alpha(\beta)\in R'$。
\end{SGAresult}

\SGAsection{2}{两个根之间的关系}
\SGAsubsection{2.1}{成比例的根}
\begin{SGAresult}{命题}{2.1.1}
设 $\alpha$ 与 $\beta$ 是两个\Term{root}{根}。下列条件\Term{equivalence}{等价}：

(i) 存在 $k\in\mathbb Q$ 使 $\alpha=k\beta$。

(ii) $s_\alpha=s_\beta$。

此外，在这些条件下，$\alpha^*=k^{-1}\beta^*$，且 $k$ 等于 $1,-1,2,-2,1/2,-1/2$ 中的一个。
\end{SGAresult}
\begin{proof}
先假设 (i)。首先有
\[\alpha^*=\ell(\alpha)^{-1}2p(\alpha)=k^{-2}\ell(\beta)^{-1}2kp(\beta)=k^{-1}\beta^*.\]
立即得到 $s_\alpha=s_\beta$。反之，若 $s_\alpha=s_\beta$，则
\[\alpha=s_\alpha(-\alpha)=s_\beta(-\alpha)=(\beta^*,\alpha)\beta-\alpha,\]
故 $2\alpha=(\beta^*,\alpha)\beta$，其中 $(\beta^*,\alpha)\in\mathbb Z$，因而 (ii) 蕴涵 (i)。最后，若 $\alpha=k\beta$，则 $\alpha^*=k^{-1}\beta^*$，从而
\[(\alpha^*,\beta)=2k^{-1},\qquad(\beta^*,\alpha)=2k,\]
所以 $2k$ 与 $2k^{-1}$ 都是整数，证明完毕。
\end{proof}

\SourcePage{68}
\begin{SGAresult}{应用}{2.1.2}
满足 $\operatorname{rgss}(\mathscr R)=1$ 的\Term{root-datum}{根数据} $\mathscr R$ 属于下列两种类型之一：

(i) $A_1$ 型：存在 $\alpha\in M$，使\Term{root}{根}为 $\alpha,-\alpha$。此时\Term{coroot}{余根}为 $\alpha^*,-\alpha^*$。

(ii) $A'_1$ 型：存在 $\alpha\in M$，使\Term{root}{根}为 $\alpha,-\alpha,2\alpha,-2\alpha$。此时\Term{coroot}{余根}为 $\alpha^*,-\alpha^*,\alpha^*/2,-\alpha^*/2$。
\end{SGAresult}
\begin{SGAresult}{定义}{2.1.3}
若 $\alpha/2\notin R$，称 $\alpha\in R$ \Term{indivisible-root}{不可分}。若每个\Term{root}{根}都\Term{indivisible-root}{不可分}，称 $\mathscr R$ 是\Term{reduced-root-datum}{约化}的。

$\mathscr R$ \Term{reduced-root-datum}{约化}的充要条件是 $\mathscr R^*$ \Term{reduced-root-datum}{约化}。若 $\alpha$ \Term{indivisible-root}{不可分}且 $w\in W$，则 $w(\alpha)$ 也\Term{indivisible-root}{不可分}。
\end{SGAresult}
\begin{SGAresult}{定义}{2.1.4}
设 $\alpha\in R$。若 $\alpha$ \Term{indivisible-root}{不可分}，置 $\operatorname{ind}(\alpha)=\alpha$。否则置 $\operatorname{ind}(\alpha)=\alpha/2$。
\end{SGAresult}
\begin{SGAresult}{推论}{2.1.5}
若 $\alpha\in R$ \Term{indivisible-root}{不可分}，且 $k\alpha\in R$，其中 $k\in\mathbb Q$，则 $k\in\mathbb Z$。
\end{SGAresult}
\begin{SGAresult}{命题}{2.1.6}
设 $\mathscr R$ 为一个\Term{root-datum}{根数据}。则
\[\operatorname{ind}(\mathscr R)=(M,M^*,\operatorname{ind}(R),\operatorname{ind}(R)^*)\]
是一个\Term{reduced-root-datum}{约化}\Term{root-datum}{根数据}，并且
\[W(\operatorname{ind}(\mathscr R))=W(\mathscr R).\]
\end{SGAresult}
\begin{proof}
事实上，\Term{weyl-group}{Weyl 群}置换\Term{indivisible-root}{不可分}\Term{root}{根}，故由 1.2.11，$\operatorname{ind}(\mathscr R)$ 为\Term{root-datum}{根数据}。第二个断言由 2.1.1 得出。
\end{proof}
\begin{SGAresult}{注记}{2.1.7}
若 $\mathscr R$ 非\Term{reduced-root-datum}{约化}，则 $\operatorname{ind}(R)^*\ne\operatorname{ind}(R^*)$，因而 $\operatorname{ind}(\mathscr R^*)\ne\operatorname{ind}(\mathscr R)^*$。
\end{SGAresult}
\SGAsubsection{2.2}{正交根}
\begin{SGAresult}{引理}{2.2.1}
设 $\alpha$ 与 $\beta$ 是两个\Term{root}{根}。有
\[\ell(\alpha)(\alpha^*,\beta)=\ell(\beta)(\beta^*,\alpha).\tag{11}\]
\end{SGAresult}
这由 1.2.1 的公式 (9) 与 (10) 立即得出。
\begin{SGAresult}{推论}{2.2.2}
设 $\alpha,\beta\in R$。下列条件\Term{equivalence}{等价}：
\begin{align*}
\text{(i)}\quad &(\alpha^*,\beta)=0,&\text{(i bis)}\quad &(\beta^*,\alpha)=0,\\
\text{(ii)}\quad &(p(\alpha),\beta)=0,\\
\text{(iii)}\quad &s_\alpha(\beta)=\beta,&\text{(iii bis)}\quad &s_\beta(\alpha)=\alpha,\\
\text{(iv)}\quad &s_\alpha(\beta^*)=\beta^*,&\text{(iv bis)}\quad &s_\beta(\alpha^*)=\alpha^*,\\
\text{(v)}\quad &s_\alpha\ne s_\beta\text{，且 }s_\alpha\text{ 与 }s_\beta\text{ 可交换。}
\end{align*}
\end{SGAresult}
\begin{proof}
除涉及 (v) 的\Term{equivalence-relation}{等价关系}外，其他都显然成立。先证 (i)（以及 (i bis)）蕴涵 (v)。有
\[s_\alpha s_\beta(x)=x-(\beta^*,x)\beta-(\alpha^*,x)\alpha+(\beta^*,x)(\alpha^*,\beta)\alpha.\]
若 $(\alpha^*,\beta)=0$，则 $(\beta^*,\alpha)=0$，且 $s_\alpha s_\beta(x)=s_\beta s_\alpha(x)$。
\SourcePage{69}
反之，假设 (v)。此时
\[s_\alpha=s_\beta s_\alpha s_\beta=s_{s_\beta(\alpha)}\qquad\text{（由 1.2.10）。}\]
由 2.1，存在 $k\in\mathbb Q^*$ 使 $\alpha=ks_\beta(\alpha)=k\alpha-k(\beta^*,\alpha)\beta$。因 $s_\alpha\ne s_\beta$，由 2.1.1，$\beta$ 与 $\alpha$ 不成比例，故 $(\beta^*,\alpha)=0$。
\end{proof}
\begin{SGAresult}{定义}{2.2.3}
满足 2.2.2 中\Term{equivalence}{等价}条件的两个\Term{root}{根}称为\Term{orthogonal-roots}{正交}的。
\end{SGAresult}
\begin{SGAresult}{注记}{2.2.4}
\Term{root}{根} $\alpha$ 与 $\beta$ \Term{orthogonal-roots}{正交}，当且仅当\Term{coroot}{余根} $\alpha^*$ 与 $\beta^*$ \Term{orthogonal-roots}{正交}。
\end{SGAresult}
\begin{SGAresult}{引理}{2.2.5}
若 $\alpha$ 与 $\beta$ 是两个\Term{orthogonal-roots}{正交}\Term{root}{根}，则 $\alpha+\beta\in R$ 当且仅当 $\alpha-\beta\in R$。
\end{SGAresult}
事实上，$s_\beta(\alpha+\beta)=s_\beta(\alpha)+s_\beta(\beta)=\alpha-\beta$。
\begin{SGAresult}{引理}{2.2.6}
设 $\alpha$ 与 $\beta$ 是两个不\Term{orthogonal-roots}{正交}的\Term{root}{根}。若对\Term{coroot}{余根} $\gamma^*$，把 $\ell(\gamma^*)$ 定义为 $\mathscr R^*$ 的\Term{root}{根} $\gamma^*$ 的 $\ell$ 值，则有关系
\[\ell(\alpha)\ell(\alpha^*)=\ell(\beta)\ell(\beta^*).\tag{12}\]
\end{SGAresult}
\begin{proof}
事实上，将公式 (11) 乘以对 $\mathscr R^*$ 而言的相应公式，并注意对每个 $\gamma\in R$ 有 $(\gamma^*)^*=\gamma$，得到
\[(\beta^*,\alpha)(\alpha^*,\beta)\ell(\alpha)\ell(\alpha^*)=(\beta^*,\alpha)(\alpha^*,\beta)\ell(\beta)\ell(\beta^*).\]
\end{proof}
\SGAsubsection{2.3}{一般情形}
\begin{SGAresult}{命题}{2.3.1}
若 $\alpha$ 与 $\beta$ 是任意两个\Term{root}{根}，则
\[0\leq(\alpha^*,\beta)(\beta^*,\alpha)\leq4.\tag{13}\]
若 $\alpha$ 与 $\beta$ 既不成比例，也不\Term{orthogonal-roots}{正交}，则
\[1\leq(\alpha^*,\beta)(\beta^*,\alpha)\leq3.\]
\end{SGAresult}
\begin{proof}
事实上，$4(p(\alpha),\beta)^2=\ell(\alpha)\ell(\beta)(\alpha^*,\beta)(\beta^*,\alpha)$。另一方面，由 1.2.6，对称\Term{bilinear}{双线性}\Term{form}{形式} $(p(x),y)$ 在 $\mathscr V(R)$ 上正定，故 $(p(x),y)^2\leq\ell(x)\ell(y)$。\footnote{编者注：等号成立当且仅当 $x$ 与 $y$ 成比例。}
\end{proof}
\begin{SGAresult}{推论}{2.3.2}
设 $\alpha$ 与 $\beta$ 是两个不\Term{orthogonal-roots}{正交}的\Term{root}{根}。若 $\ell(\alpha)=\ell(\beta)$，则存在 $w\in W$ 使 $w(\alpha)=\beta$。
\end{SGAresult}
\begin{proof}
事实上，若 $\alpha$ 与 $\beta$ 成比例，则由 $\ell(\alpha)=\ell(\beta)$ 得 $\alpha=\beta$ 或 $\alpha=-\beta$；此时取 $w=1$ 或 $w=s_\alpha$。若 $\alpha$ 与 $\beta$ 既不成比例，也不\Term{orthogonal-roots}{正交}，由公式 (11) 与 2.3.1，
\[(\beta^*,\alpha)=(\alpha^*,\beta)=\pm1.\]
若 $(\beta^*,\alpha)=(\alpha^*,\beta)=1$，取 $w=s_\beta s_\alpha s_\beta$；若 $(\beta^*,\alpha)=(\alpha^*,\beta)=-1$，取 $w=s_\alpha s_\beta$。
\end{proof}

\SourcePage{70}
\begin{SGAresult}{推论}{2.3.3}
若 $\alpha$ 与 $\beta$ 是两个\Term{root}{根}，且 $\alpha\ne\beta$、$(\beta^*,\alpha)>0$（分别地，$\alpha\ne-\beta$、$(\beta^*,\alpha)<0$），则 $\alpha-\beta$（分别地，$\alpha+\beta$）是一个\Term{root}{根}。
\end{SGAresult}
\begin{proof}
将 $\beta$ 换成 $-\beta$，即可由第一种情形得到第二种。若 $\alpha$ 与 $\beta$ 成比例，且 $(\beta^*,\alpha)>0$，则 $\beta=\alpha$、$2\beta=\alpha$ 或 $\beta=2\alpha$。第一种可能已排除。在其余两种情形中，分别有 $\alpha-\beta=\beta\in R$ 或 $\alpha-\beta=-\alpha\in R$。若 $\alpha$ 与 $\beta$ 不成比例，则 $(\alpha^*,\beta)$ 与 $(\beta^*,\alpha)$ 是两个严格正的整数，乘积至多为 $3$，故其中一个等于 $1$。若 $(\beta^*,\alpha)=1$，则 $\alpha-\beta=s_\beta(\alpha)\in R$；若 $(\alpha^*,\beta)=1$，则 $\alpha-\beta=-s_\alpha(\beta)\in R$。
\end{proof}
\begin{SGAresult}{引理}{2.3.4}
设 $\alpha$ 与 $\beta$ 是两个不成比例的\Term{root}{根}。若 $\beta-\alpha$ 不是\Term{root}{根}，则当 $k=-(\alpha^*,\beta)$ 时 $\beta+k\alpha\in R$，但当 $k=-(\alpha^*,\beta)+1$ 时则不然。
\end{SGAresult}
事实上，$\beta-(\alpha^*,\beta)\alpha=s_\alpha(\beta)\in R$，而 $\beta+\bigl(-(\alpha^*,\beta)+1\bigr)\alpha=s_\alpha(\beta-\alpha)\notin R$。
\begin{SGAresult}{命题}{2.3.5}
设 $\alpha$ 与 $\beta$ 是两个不成比例的\Term{root}{根}。使 $\beta+k\alpha\in R$ 的整数 $k$ 的\Term{set}{集合}是一个区间 $[p,q]$（$p\leq0$，$q\geq0$），且 $p+q=-(\alpha^*,\beta)$。
\end{SGAresult}
\begin{proof}
对于第一个断言，例如只须证明：若 $\beta+k\alpha\in R$，其中整数 $k>0$，则 $\beta+(k-1)\alpha\in R$。当 $k=1$ 时显然成立。当 $k\geq2$ 时，
\[(\alpha^*,\beta+k\alpha)=(\alpha^*,\beta)+2k\geq-3+4>0,\]
再用 2.3.3 即可。于是设所论区间为 $[p,q]$。对 $\beta+p\alpha$ 应用 2.3.4，得到
\[q-p=-(\alpha^*,\beta+p\alpha)=-(\alpha^*,\beta)-2p.\]
\footnote{编者注：已将 $+2p$ 改正为 $-2p$。}
\end{proof}
\begin{SGAresult}{注记}{2.3.6}
上述公式包含定性断言 2.2.5 与 2.3.3。
\end{SGAresult}
\begin{SGAresult}{补充}{2.3.7}
\footnote{编者注：增加了这些补充，它们对 3.1.5 的证明有用。}由 1.2.1 (9) 与 1.2.6，$\mathscr V(R)$ 上由
\[\langle x,y\rangle=(p(x),y)\]
定义的\Term{bilinear}{双线性}\Term{form}{形式}是对称正定的。由 1.2.1 (10)，对每个 $\alpha\in R$ 与 $y\in\mathscr V(R)$，有
\[\langle\alpha,y\rangle=\frac{\ell(\alpha)}2(\alpha^*,y),\]
其中 $\ell(\alpha)=\langle\alpha,\alpha\rangle$。因此，由 2.3.3 得到下列推论。
\end{SGAresult}
\begin{SGAresult}{推论}{}
设 $\alpha\ne\beta$ 是 $R$ 中的元素。若 $\alpha-\beta\notin R$，则 $\langle\alpha,\beta\rangle\leq0$。
\end{SGAresult}

\SourcePage{71}
\SGAsection{3}{单根与正根}
\SGAsubsection{3.1}{单根系}
\begin{SGAresult}{引理}{3.1.1}
设 $\alpha$ 与 $\alpha_i$（$i=1,\ldots,n$）是\Term{root}{根}。假设 $\alpha$ 与各 $\alpha_i$ 均不同。若有关系
\[\alpha=\sum_{i=1}^nq_i\alpha_i,\qquad q_i\in\mathbb Q_+,\]
则存在一个指标 $i$，使 $q_i\ne0$、$(\alpha^*,\alpha_i)>0$、$\alpha-\alpha_i\in R$。
\end{SGAresult}
\begin{proof}
事实上，写出 $2=(\alpha^*,\alpha)=\sum_{i=1}^nq_i(\alpha^*,\alpha_i)$，便证明了前两个断言。第三个断言由 2.3.3 得出。
\end{proof}
\begin{SGAresult}{命题}{3.1.2}
设 $\alpha$ 与 $\alpha_i$（$i=1,\ldots,n$）是\Term{root}{根}。若
\[\alpha=\sum_{i=1}^nm_i\alpha_i,\qquad m_i\in\mathbb N_+,\]
则存在由各 $\alpha_i$ 中选出的\Term{root}{根}的序列 $\beta_1,\ldots,\beta_m$（不必两两不同），使得若记
\[\gamma_p=\sum_{i=1}^p\beta_i,\qquad p=1,\ldots,m,\]
则 $\gamma_p\in R$ 且 $\gamma_m=\alpha$。
\end{SGAresult}
\begin{proof}
对整数 $\sum m_i=m'$ 归纳。若 $\alpha$ 等于某个 $\alpha_i$，设为 $\alpha_{i_0}$（当 $m'=1$ 时必然如此），则取 $m=1$、$\beta_1=\alpha_{i_0}$。否则应用前一引理，可得指标 $i$ 使 $m_i\ne0$ 且 $\alpha-\alpha_i$ 是\Term{root}{根}。此时 $(m_i-1)\in\mathbb N$，并且
\[\alpha-\alpha_i=(m_i-1)\alpha_i+\sum_{j\ne i}m_j\alpha_j.\]
对 $\alpha-\alpha_i$ 应用归纳假设即可。
\end{proof}
\begin{SGAresult}{推论}{3.1.3}
设 $R'\subset R$。下列条件\Term{equivalence}{等价}：

(i) 若 $\alpha,\beta\in R'$ 且 $\alpha+\beta\in R$，则 $\alpha+\beta\in R'$。

(ii) $(\mathbb N\cdot R')\cap R\subset R'$。
\end{SGAresult}
事实上，显然 (ii) $\Rightarrow$ (i)。反向蕴涵由上述命题立即得出。
\begin{SGAresult}{定义}{3.1.4}
满足 3.1.3 中条件的\Term{root}{根}的\Term{set}{集合}称为\Term{closed-root-set}{闭}的。
\end{SGAresult}
\begin{SGAresult}{命题}{3.1.5}
设 $\Delta\subset R$ 是一个\Term{root}{根}的\Term{set}{集合}。下列条件\Term{equivalence}{等价}：

(i) $\Delta$ 的元素\Term{indivisible-root}{不可分}、\Term{linearly-independent}{线性无关}，且
\[R\subset(\mathbb Q_+\cdot\Delta)\cup(-\mathbb Q_+\cdot\Delta).\]
\SourcePage{72}
(ii) $\Delta$ 的元素\Term{linearly-independent}{线性无关}，且
\[R\subset(\mathbb N\cdot\Delta)\cup(-\mathbb N\cdot\Delta).\]
(iii) 每个\Term{root}{根}唯一地写成 $\Delta$ 中元素的线性组合，其系数是同号的整数。
\end{SGAresult}
\begin{proof}
\footnote{编者注：下文调整了顺序，并详细展开了 (iii) $\Rightarrow$ (i) 的证明。}显然有 (ii) $\Rightarrow$ (iii)。以 $\alpha_1,\ldots,\alpha_n$ 表示 $\Delta$ 的各个（不同的）元素。

证明 (i) $\Rightarrow$ (ii)。设 $\alpha\in R$。于是有唯一的表示
\[\varepsilon\alpha=\sum q_i\alpha_i,\qquad q_i\in\mathbb Q_+,\quad\varepsilon=\pm1.\]
证明各 $q_i$ 均为整数。若除一个外所有系数都为零，这当然成立（见 2.1.5）。否则 $\alpha$ 与各 $\alpha_i$ 不同；应用 3.1.1，得到某个 $i_0$ 使 $\alpha'=\varepsilon\alpha-\alpha_{i_0}\in R$ 且 $q_{i_0}\ne0$。于是
\[\alpha'=(q_{i_0}-1)\alpha_{i_0}+\sum_{i\ne i_0}q_i\alpha_i.\]
由于至少一个 $q_i$（$i\ne i_0$）不为零，由 (i) 得 $q_{i_0}-1\geq0$。对 $\alpha'$ 重复这一操作，经过\Term{finite}{有限}步便证明各 $q_i$ 都是整数。

证明 (iii) $\Rightarrow$ (i)。先证 $\alpha_1,\alpha_2,\ldots,\alpha_n$ 在 $\mathbb Q$ 上\Term{linearly-independent}{线性无关}。否则存在等式
\[x=\sum_{i\in I}a_i\alpha_i=\sum_{j\in J}b_j\alpha_j,\]
其中 $I,J$ 是 $\{1,\ldots,n\}$ 的两个\Term{disjoint}{不相交}的\Term{subset}{子集}，至少一个（设为 $I$）\Term{nonempty}{非空}，且 $a_i,b_j\in\mathbb N^*$。由推论 2.3.7，当 $i\in I$、$j\in J$ 时，$\langle\alpha_i,\alpha_j\rangle\leq0$，故 $\langle x,x\rangle\leq0$，从而 $x=0$。取 $i_0\in I$，则等式
\[\alpha_{i_0}=\alpha_{i_0}+\sum_{i\in I}a_i\alpha_i
=\alpha_{i_0}+\sum_{j\in J}b_j\alpha_j\]
由 (iii) 蕴涵对所有 $i,j$ 都有 $a_i=0=b_j$，矛盾。这就证明 $\Delta$ 的元素在 $\mathbb Q$ 上\Term{linearly-independent}{线性无关}；再证明它们也\Term{indivisible-root}{不可分}。

设 $\beta\in\Delta$ \Term{separable}{可分}。则 $\beta/2\in R$，从而
\[\frac\beta2=\varepsilon\sum_{\alpha\in\Delta}m_\alpha\alpha,\qquad m_\alpha\in\mathbb N,\quad\varepsilon=\pm1,\]
所以也有 $\beta=2\varepsilon\sum_\alpha m_\alpha\alpha$；由唯一性，$\alpha\ne\beta$ 时 $m_\alpha=0$，且 $2\varepsilon m_\beta=1$，矛盾。
\end{proof}
\begin{SGAresult}{定义}{3.1.6}
满足 3.1.5 中条件的\Term{root}{根}的\Term{set}{集合} $\Delta$，称为一个\Term{simple-root-system}{单根系}，或 $R$ 的一个\Term{root-basis}{基}。

若 $w\in W$ 且 $\Delta$ 是一个\Term{simple-root-system}{单根系}，则 $w(\Delta)$ 也是一个\Term{simple-root-system}{单根系}。
\end{SGAresult}

\SourcePage{73}
\begin{SGAresult}{注记}{3.1.7}
这个定义只涉及 $R$ 而不涉及 $\mathscr R$；事实上，它甚至只涉及 $\operatorname{ind}(R)$。
\end{SGAresult}
\begin{SGAresult}{注记}{3.1.8}
若 $\Delta$ 是一个\Term{simple-root-system}{单根系}，则 $\Delta$ 是自由\Term{abelian-group}{阿贝尔群} $\Gamma_0(R)$ 的一组基。故 $\operatorname{Card}(\Delta)=\operatorname{rgss}(\mathscr R)$。
\end{SGAresult}
\begin{SGAresult}{注记}{3.1.9}
于是 3.1.5 中的条件显然\Term{equivalence}{等价}于下列条件：

(i$'$) $\Delta$ 的元素\Term{indivisible-root}{不可分}，个数 $\leq\operatorname{rgss}(R)$，且 $R\subset(\mathbb Q_+\cdot\Delta)\cup(-\mathbb Q_+\cdot\Delta)$。

(ii$'$) $\operatorname{Card}(\Delta)\leq\operatorname{rgss}(R)$，且 $R\subset(\mathbb N\cdot\Delta)\cup(-\mathbb N\cdot\Delta)$。
\end{SGAresult}
\begin{SGAresult}{推论}{3.1.10}
若 $\Delta$ 是一个\Term{simple-root-system}{单根系}，则 $\operatorname{ind}(\Delta^*)$ 是一个\Term{simple-coroot-system}{单余根系}（即 $R^*$ 的一个\Term{simple-root-system}{单根系}）。
\end{SGAresult}
\begin{proof}
事实上，若 $\beta=\sum_{\alpha\in\Delta}a_\alpha\alpha$（$a_\alpha\in\mathbb N$），则 $p(\alpha)=\sum_{\alpha\in\Delta}a_\alpha p(\alpha)$，从而由 1.2.1 (10)，
\[\beta^*=\sum_{\alpha\in\Delta}a_\alpha\frac{\ell(\alpha)}{\ell(\beta)}\alpha^*,\]
这表明 $\operatorname{ind}(\Delta^*)$ 满足 (i$'$)。
\end{proof}
由 2.1.1，若 $\alpha^*$ 不\Term{indivisible-root}{不可分}，则 $\alpha^*=2u^*$，其中 $u^*\in\operatorname{ind}(\Delta^*)$（且 $u=2\alpha$）。\footnote{编者注：增加了这句话；并在 3.1.11 中将 $4a_\alpha\ell(\alpha)$ 改为 $2a_\alpha\ell(\alpha)$。}由此得到：
\begin{SGAresult}{推论}{3.1.11}
若 $\Delta$ 是一个\Term{simple-root-system}{单根系}，且 $\beta=\sum_{\alpha\in\Delta}a_\alpha\alpha$（$a_\alpha\in\mathbb Z$）是 $\beta$ 关于 $\Delta$ 的表示，则 $2a_\alpha\ell(\alpha)$ 可被 $\ell(\beta)$ 整除（若 $\alpha^*$ \Term{indivisible-root}{不可分}，甚至可被 $2\ell(\beta)$ 整除）。
\end{SGAresult}
在继续陈述\Term{simple-root-system}{单根系}的性质之前，先证明它们存在。

\SGAsubsection{3.2}{正根系}
\begin{SGAresult}{定义}{3.2.1}
若\Term{set}{集合} $R_+\subset R$ 满足下列条件，就称其为 $\mathscr R$（或 $R$，见注记 3.2.2）的一个\Term{positive-root-system}{正根系}：
\begin{align*}
\text{(P 1)}&\quad R_+\cap-R_+=\varnothing.\\
\text{(P 2)}&\quad R_+\cup-R_+=R.\\
\text{(P 3)}&\quad(\mathbb Q_+\cdot R_+)\cap R\subset R_+.
\end{align*}
特别地，这样的\Term{set}{集合}是\Term{closed-root-set}{闭}的。后文将看到，实际上，满足 (P 1) 与 (P 2) 的\Term{closed-root-set}{闭}集也满足 (P 3)，因而是一个\Term{positive-root-system}{正根系}。若 $w\in W$ 且 $R_+$ 是一个\Term{positive-root-system}{正根系}，则 $w(R_+)$ 也是一个\Term{positive-root-system}{正根系}。
\end{SGAresult}
\begin{SGAresult}{注记}{3.2.2}
这个定义只涉及 $R$。也说 $R_+$ 是 $R$ 的一个\Term{positive-root-system}{正根系}。
\end{SGAresult}

\SourcePage{74}
\begin{SGAresult}{注记}{3.2.3}
由 (P 1) 与 (P 2) 立即得到
\[\operatorname{Card}(R_+)=\operatorname{Card}(R)/2.\]
因此，若 $R_+$ 与 $R'_+$ 是两个\Term{positive-root-system}{正根系}，且 $R_+\subset R'_+$，则 $R_+=R'_+$。
\end{SGAresult}
\begin{SGAresult}{注记}{3.2.4}
若 $R_+$ 是一个\Term{positive-root-system}{正根系}，则 $R_+^*$ 是一个\Term{positive-coroot-system}{正余根系}（即 $R^*$ 的一个\Term{positive-root-system}{正根系}）。
\end{SGAresult}
这由 1.1.4 与 1.2.2 立即得出。
\begin{SGAresult}{定义}{3.2.5}
设 $\Delta$ 是一个\Term{simple-root-system}{单根系}。置
\[\mathscr P(\Delta)=(\mathbb Q_+\cdot\Delta)\cap R=(\mathbb N_+\cdot\Delta)\cap R.\]
\end{SGAresult}
\begin{SGAresult}{命题}{3.2.6}
若 $\Delta$ 是一个\Term{simple-root-system}{单根系}，则 $\mathscr P(\Delta)$ 是一个\Term{positive-root-system}{正根系}。若 $\Delta$ 是一个\Term{simple-root-system}{单根系}，$R_+$ 是一个\Term{positive-root-system}{正根系}，则有\Term{equivalence-relation}{等价关系}
\[\Delta\subset R_+\quad\Longleftrightarrow\quad R_+=\mathscr P(\Delta).\]
\end{SGAresult}
\begin{proof}
第一个断言显然。若 $\Delta\subset R_+$，则由 (P 3)，$\mathscr P(\Delta)\subset R_+$，故由 3.2.3，$\mathscr P(\Delta)=R_+$。其余显然。
\end{proof}
\begin{SGAresult}{注记}{3.2.7}
\Term{positive-root-system}{正根系}存在：设 $\geq$ 是 $\mathscr V(R)$ 上的一个全序\Term{vector-space}{向量空间}结构。关于这个\Term{order-relation}{序关系}的 $\geq0$ 的\Term{root}{根}的\Term{set}{集合}，就是一个\Term{positive-root-system}{正根系}。
\end{SGAresult}
\begin{SGAresult}{定理}{3.2.8}
设 $R_+$ 是一个\Term{positive-root-system}{正根系}。存在唯一的\Term{simple-root-system}{单根系} $\mathscr S(R_+)$，使 $\mathscr S(R_+)\subset R_+$，也就是说，使 $R_+=\mathscr P(\mathscr S(R_+))$。
\end{SGAresult}
唯一性由下列引理立即得出：
\begin{SGAresult}{引理}{3.2.9}
设 $\Delta$ 是一个\Term{simple-root-system}{单根系}。$\alpha\in\mathscr P(\Delta)$ 属于 $\Delta$ 的充要条件是：$\alpha$ 不能写成 $\mathscr P(\Delta)$ 中两个元素的和。
\end{SGAresult}
这个引理由定义与 3.1.2 立即得出。

现在证明 $\mathscr S(R_+)$ 的存在性。考虑 $R_+$ 的所有满足 $T=\operatorname{ind}(T)$ 且 $(\mathbb Q_+\cdot T)\supset R_+$ 的\Term{subset}{子集} $T$ 所成的\Term{set}{集合}。此\Term{set}{集合}\Term{nonempty}{非空}，因为它含有 $\operatorname{ind}(R_+)$。取其中关于包含关系的一个极小元素 $\Delta$。证明 $\Delta$ 是一个\Term{simple-root-system}{单根系}；由 3.1.5 (i)，也就是证明 $\Delta$ 是 $M\otimes\mathbb Q$ 的一个\Term{linearly-independent}{线性无关}\Term{subset}{子集}。
\begin{SGAresult}{引理}{3.2.10}
若 $\alpha,\beta\in\Delta$，且 $q\alpha+q'\beta\in R$，其中 $q,q'\in\mathbb Q$，则 $qq'\geq0$。
\end{SGAresult}
\begin{proof}
事实上，若 $qq'<0$，则（必要时交换 $\alpha$ 与 $\beta$）可写为
\[a\alpha-b\beta\in R_+,\qquad a,b\in\mathbb Q_+^*,\]
于是由假设，存在关系
\[a\alpha-b\beta=\sum_{\gamma\in S}c(\gamma)\gamma,\qquad c(\gamma)\in\mathbb Q_+.\]
\SourcePage{75}
若 $a\leq c(\alpha)$，可写成
\[-\beta=\frac{c(\alpha)-a}{b}\alpha+\sum_{\gamma\ne\alpha}\frac{c(\gamma)}b\gamma.\]
于是此元素属于 $(\mathbb Q_+\cdot\Delta)\cap R$，由 (P 3)，后者包含于 $R_+$。因而 $\beta\in R_+\cap-R_+$，与 (P 1) 矛盾。

反之，若 $a>c(\alpha)$，写成
\[(a-c(\alpha))\alpha=b\beta+\sum_{\gamma\ne\alpha}c(\gamma)\gamma,\]
这表明 $\Delta\subset\mathbb Q_+\cdot(\Delta-\{\alpha\})$，与 $\Delta$ 的极小性矛盾。
\end{proof}
回忆（见 2.3.7），$\mathscr V(R)$ 上由 $\langle x,y\rangle=(p(x),y)$ 定义的\Term{bilinear}{双线性}\Term{form}{形式}是一个欧氏内积；此外，对 $\alpha\in R$ 有 $\langle\alpha,y\rangle=\ell(\alpha)(\alpha^*,y)/2$，因此 $\langle\alpha,y\rangle$ 与 $(\alpha^*,y)$ 同号。
\begin{SGAresult}{引理}{3.2.11}
若 $\alpha,\beta\in\Delta$，则 $(\beta^*,\alpha)\leq0$，从而 $\langle\beta,\alpha\rangle\leq0$。
\end{SGAresult}
\begin{proof}
事实上，$s_\beta(\alpha)=\alpha-(\beta^*,\alpha)\beta\in R$，故由 3.2.10 得 $(\beta^*,\alpha)\leq0$。
\end{proof}
现在证明 $\Delta$ \Term{linearly-independent}{线性无关}。否则，由 3.2.11，如同在 3.1.5 的证明中，可得一个非平凡关系
\[\sum_{\alpha\in\Delta}m(\alpha)\alpha=0,\qquad m(\alpha)\in\mathbb N,\]
从而若 $m(\alpha_0)\geq1$，则 $-\alpha_0=(m(\alpha_0)-1)\alpha_0+\sum_{\alpha\ne\alpha_0}m(\alpha)\alpha$。于是，由 (P 3)，$\alpha_0$ 属于 $P\cap-P$，与 (P 1) 矛盾。

这就证明 $\Delta$ 是 $R$ 的一个\Term{root-basis}{基}，并完成了定理 3.2.8 的证明。
\begin{SGAresult}{推论}{3.2.12}
设 $R_+$ 是一个\Term{positive-root-system}{正根系}，$\Delta$ 是 $R$ 的一个\Term{root-basis}{基}，且 $w\in W$。有
\begin{gather*}
\Delta\subset R_+\quad\Longleftrightarrow\quad R_+=\mathscr P(\Delta)\quad\Longleftrightarrow\quad\Delta=\mathscr S(R_+),\\
\mathscr P(\operatorname{ind}(\Delta^*))=\mathscr P(\Delta)^*,\qquad
\mathscr S(R_+^*)=\operatorname{ind}(\mathscr S(R_+)^*);\\
\mathscr S(w(R_+))=w(\mathscr S(R_+)),\qquad\mathscr P(w(\Delta))=w(\mathscr P(\Delta)).
\end{gather*}
\end{SGAresult}
\begin{SGAresult}{定义}{3.2.13}
若已选定一个\Term{simple-root-system}{单根系} $\Delta$，则称 $\mathscr P(\Delta)$ 中的元素为\Term{positive-root}{正根}。若已选定一个\Term{positive-root-system}{正根系} $R_+$，则称 $\mathscr S(R_+)$ 中的元素为\Term{simple-root}{单根}。
\end{SGAresult}
\begin{SGAresult}{推论}{3.2.14}
设 $R_+$ 是一个\Term{positive-root-system}{正根系}，$\alpha\in R_+$。下列条件\Term{equivalence}{等价}：

(i) $\alpha$ 是\Term{simple-root}{单根}（即 $\alpha\in\mathscr S(R_+)$）。

(ii) $\alpha$ 不是 $R_+$ 中两个元素的和。

(iii) $R_+-\{\alpha\}$ 是\Term{closed-root-set}{闭}的（见 3.1.4）。
\end{SGAresult}
(i) 与 (ii) 的\Term{equivalence}{等价}性由 3.2.9 立即得出。(ii) 与 (iii) 的\Term{equivalence}{等价}性显然。

\SourcePage{76}
\begin{SGAresult}{定义}{3.2.15}
设 $\Delta$ 是一个\Term{simple-root-system}{单根系}。\Term{root}{根} $\alpha$ 关于 $\Delta$ 的\Term{resolution}{分解}中各系数之和，称为 $\alpha$ 关于 $\Delta$ 的\Term{root-order}{阶}，记为 $\operatorname{ord}_\Delta(\alpha)$。

有\Term{equivalence-relation}{等价关系}：
\[\alpha\in\mathscr P(\Delta)\Longleftrightarrow\operatorname{ord}_\Delta(\alpha)>0,\qquad
\alpha\in\Delta\Longleftrightarrow\operatorname{ord}_\Delta(\alpha)=1.\]
\end{SGAresult}
\begin{SGAresult}{引理}{3.2.16}
设 $\Delta$ 是一个\Term{simple-root-system}{单根系}，$\alpha\in\mathscr P(\Delta)$。存在一个序列 $\alpha_i\in\Delta$（$i=1,\ldots,m$），使
\begin{gather*}
\alpha_1+\alpha_2+\cdots+\alpha_p\in\mathscr P(\Delta)\qquad(p=1,\ldots,m),\\
\alpha_1+\alpha_2+\cdots+\alpha_m=\alpha.
\end{gather*}
此外，对满足这些条件的任一序列 $\alpha_i$，都有 $m=\operatorname{ord}_\Delta(\alpha)$。
\end{SGAresult}
由 3.1.2 显然成立。

\SGAsubsection{3.3}{正根系的刻画与共轭}
\begin{SGAresult}{引理}{3.3.1}
若 $\alpha\in\mathscr P(\Delta)$、$\beta\in\Delta$，且 $\alpha$ 与 $\beta$ 不成比例，则 $s_\beta(\alpha)\in\mathscr P(\Delta)$。
\end{SGAresult}
\begin{proof}
事实上，$s_\beta(\alpha)=\alpha-(\beta^*,\alpha)\beta$。由于 $\alpha$ 的\Term{resolution}{分解}中至少有一个异于 $\beta$ 的\Term{simple-root}{单根}以非零（因而严格正的）系数出现，它在 $s_\beta(\alpha)$ 的\Term{resolution}{分解}中也以相同的系数出现，所以 $s_\beta(\alpha)$ 也是\Term{positive-root}{正根}。
\end{proof}
\begin{SGAresult}{推论}{3.3.2}
若 $\beta\in\Delta$，则对称变换 $s_\beta$ 置换 $\mathscr P(\Delta)$ 中不与 $\beta$ 成比例的元素。
\end{SGAresult}
\begin{SGAresult}{引理}{3.3.3}
若 $\alpha\in\mathscr P(\Delta)-\Delta$，且 $\alpha$ \Term{indivisible-root}{不可分}，则存在 $\beta\in\Delta$，使 $s_\beta(\alpha)\in\mathscr P(\Delta)$，且 $\operatorname{ord}_\Delta(s_\beta(\alpha))<\operatorname{ord}_\Delta(\alpha)$。
\end{SGAresult}
\begin{proof}
事实上，由 3.1.1，存在 $\beta\in\Delta$ 使 $(\beta^*,\alpha)>0$。由于 $\alpha\notin\Delta$ 且\Term{indivisible-root}{不可分}，$\alpha$ 不可能与 $\beta$ 成比例。因此，由 3.3.1，$s_\beta(\alpha)\in\mathscr P(\Delta)$，且 $\operatorname{ord}_\Delta(s_\beta(\alpha))=\operatorname{ord}_\Delta(\alpha)-(\beta^*,\alpha)<\operatorname{ord}_\Delta(\alpha)$。
\end{proof}
\begin{SGAresult}{推论}{3.3.4}
若 $\alpha\in\mathscr P(\Delta)$ \Term{indivisible-root}{不可分}，则存在一个序列 $\beta_p\in\Delta$（$p=1,\ldots,q$），使
\[\alpha_p=s_{\beta_p}\cdots s_{\beta_1}(\alpha)\in\mathscr P(\Delta)\qquad(p=1,\ldots,q),\]
且 $\alpha_q\in\Delta$。
\end{SGAresult}
对 $\operatorname{ord}_\Delta(\alpha)$ 归纳，由 3.3.3 即得。
\begin{SGAresult}{命题}{3.3.5}
\Term{weyl-group}{Weyl 群}由 $s_\alpha$（$\alpha\in\Delta$）生成。每个\Term{indivisible-root}{不可分}\Term{root}{根}都在\Term{weyl-group}{Weyl 群}某个元素\Term{action}{作用}下与一个\Term{simple-root}{单根}\Term{conjugacy}{共轭}。
\end{SGAresult}
\begin{proof}
第二个断言由 3.3.4 立即得出。再由 1.2.10 与 2.1.1 得到第一个断言。
\end{proof}
\begin{SGAresult}{命题}{3.3.6}
设 $R_+$ 是一个\Term{positive-root-system}{正根系}。设 $P'\subset R$ 满足 (P 2) 且为\Term{closed-root-set}{闭}集。则存在 $w\in W$ 使 $w(R_+)\subset P'$。
\end{SGAresult}

\SourcePage{77}
先立即陈述其推论。
\begin{SGAresult}{推论}{3.3.7}
\Term{weyl-group}{Weyl 群}在全体\Term{positive-root-system}{正根系}（分别地，全体\Term{simple-root-system}{单根系}）的\Term{set}{集合}上\Term{transitive-action}{传递地作用}。
\end{SGAresult}
\begin{SGAresult}{推论}{3.3.8}
$R$ 的一个\Term{subset}{子集}为\Term{positive-root-system}{正根系}，当且仅当它满足 (P 1)、(P 2)，且为\Term{closed-root-set}{闭}集。
\end{SGAresult}
\begin{SGAresult}{推论}{3.3.9}
若在 $\Gamma_0(R)$ 上赋予一个有序\Term{group-structure}{群结构}，使每个\Term{root}{根}都 $>0$ 或 $<0$，则关于这个序结构的\Term{positive-root}{正根}的\Term{set}{集合}是一个\Term{positive-root-system}{正根系}。
\end{SGAresult}
\begin{proof}[3.3.6 的证明]
可取 $w\in W$，使 $\operatorname{Card}(w(R_+)\cap P')$ 最大。因此，将 $R_+$ 换成 $w(R_+)$，可假设对每个 $\alpha\in\mathscr S(R_+)=\Delta$ 都有
\[\operatorname{Card}(R_+\cap P')\geq\operatorname{Card}(R_+\cap s_\alpha(P')).\tag{*}\]
证明 $R_+\subset P'$。否则，由于 $P'$ 是\Term{closed-root-set}{闭}集，存在 $\alpha\in\Delta$，$\alpha\notin P'$。但 $P'$ 满足 (P 2)，故 $-\alpha\in P'$（因此，若 $2\alpha$ 是\Term{root}{根}，也有 $-2\alpha\in P'$）。对每个 $\beta\in R_+\cap P'$，有 $\beta\ne\alpha$；若 $2\alpha$ 不是\Term{root}{根}，则由 3.3.1 得 $s_\alpha(\beta)\in R_+$，因而
\[s_\alpha(R_+\cap P')\subset R_+\cap s_\alpha(P');\]
但 $R_+\cap s_\alpha(P')$ 还包含 $\alpha=s_\alpha(-\alpha)$，这与不等式 (*) 矛盾。若 $2\alpha$ 是\Term{root}{根}，同理。
\end{proof}
因此，对满足 (P 2) 且为\Term{closed-root-set}{闭}集的\Term{root}{根}的\Term{set}{集合}的研究，归结为它包含一个\Term{positive-root-system}{正根系}的情形。
\begin{SGAresult}{命题}{3.3.10}
设 $R_+$ 是一个\Term{positive-root-system}{正根系}，$P'$ 是 $R$ 的包含 $R_+$ 的一个\Term{closed-root-set}{闭}\Term{subset}{子集}。记 $\Delta=\mathscr S(R_+)$，$\Delta'=\Delta\cap-P'$，则 $P'$ 是 $R_+$ 与那些可写成 $\Delta'$ 中元素的负系数线性组合的\Term{root}{根}的\Term{set}{集合}之并。
\end{SGAresult}
\begin{proof}
对\Term{root}{根} $\gamma\in P'-R_+$ 的\Term{root-order}{阶}归纳证明。若 $\operatorname{ord}_\Delta(\gamma)=-1$，则 $-\gamma\in\Delta$，且 $\gamma\in-\Delta'$。若 $\operatorname{ord}_\Delta(\gamma)<-1$，由 3.1.2，存在 $\beta\in\Delta$ 使 $-\gamma-\beta\in R$；于是
\[0<\operatorname{ord}_\Delta(-\gamma-\beta)=\operatorname{ord}_\Delta(-\gamma)-1<\operatorname{ord}_\Delta(-\gamma),\]
第一个不等式表明 $-\gamma-\beta\in R_+$。所以，由于 $R_+\cap-R_+=\varnothing$，且 $\gamma+\beta$ 是 $P'$ 中两个\Term{root}{根}的和，它是 $P'-R_+$ 的一个元素，且 $\operatorname{ord}_\Delta(\gamma+\beta)>\operatorname{ord}_\Delta(\gamma)$。因此，由归纳假设，$\gamma+\beta$ 是 $\Delta'$ 中元素的负系数线性组合。因 $\gamma=(\gamma+\beta)-\beta$，只须验证 $\beta\in-P'$。而 $\gamma\in P'$，$-(\gamma+\beta)\in R_+\subset P'$，所以 $-\beta=\gamma-(\gamma+\beta)$ 属于 $P'$。
\end{proof}
\begin{SGAresult}{定义}{3.3.10.1}
\footnote{编者注：在此插入了这个定义。}若 $R$ 的\Term{subset}{子集} $R'$ 满足 $-R'=R'$，则称其为\Term{symmetric-root-set}{对称}的。
\end{SGAresult}

\SourcePage{78}
\begin{SGAresult}{说明}{3.3.11}
设 $P'$ 是一个满足 (P 2) 且\Term{closed-root-set}{闭}的\Term{root}{根}的\Term{set}{集合}。存在一个\Term{simple-root-system}{单根系} $\Delta$ 与 $\Delta$ 的一个\Term{subset}{子集} $\Delta'$，使
\[P'=R\cap(\mathbb N\cdot\Delta\cup-\mathbb N\cdot\Delta').\]
若记 $R_{\Delta'}=(\mathbb Z\cdot\Delta')\cap R$，这是 $R$ 的一个\Term{closed-root-set}{闭}且\Term{symmetric-root-set}{对称}的\Term{subset}{子集}，则 $P'$ 是 $R_{\Delta'}$ 与 $\mathscr P(\Delta)$ 的如下\Term{closed-root-set}{闭}\Term{subset}{子集}的\Term{disjoint-union}{不交并}：该\Term{subset}{子集}由 $\alpha\in\mathscr P(\Delta)-\mathbb N\cdot\Delta'$ 构成，也就是说，由那些在关于 $\Delta$ 的\Term{resolution}{分解}中“至少包含一个 $\Delta-\Delta'$ 中的\Term{root}{根}”的\Term{positive-root}{正根}构成。
\end{SGAresult}

\SGAsubsection{3.4}{闭且对称的根集}
\begin{SGAresult}{命题}{3.4.1}
设 $\mathscr R=(M,M^*,R,R^*)$ 是一个\Term{root-datum}{根数据}，$R'$ 是 $R$ 的一个\Term{closed-root-set}{闭}且\Term{symmetric-root-set}{对称}的\Term{subset}{子集}。则：

(i) $\mathscr R'=(M,M^*,R',R'^*)$ 是一个\Term{root-datum}{根数据}；

(ii) 对 $R$ 的任一\Term{positive-root-system}{正根系} $R_+$，$R'_+=R_+\cap R'$ 是 $R'$ 的一个\Term{positive-root-system}{正根系}；

(iii) $\mathscr R'$ 的\Term{weyl-group}{Weyl 群} $W(\mathscr R')$ 是 $W(\mathscr R)$ 中由 $s_\alpha$（$\alpha\in R'$）生成的\Term{subgroup}{子群}。
\end{SGAresult}
\begin{proof}
第一个断言由 1.2.11 显然成立；第二个由 3.3.8 得出；第三个显然。
\end{proof}
\begin{SGAresult}{推论}{3.4.2}
设 $w\in W(\mathscr R')$。$w$ 的阶是使每个 $\alpha'\in R'$ 都满足 $w^n(\alpha')=\alpha'$ 的最小整数 $n>0$。
\end{SGAresult}
只须对\Term{root-datum}{根数据} $\mathscr R'$ 应用 1.2.7。
\begin{SGAresult}{推论}{3.4.3}
设 $\alpha$ 与 $\beta$ 是两个不成比例的\Term{root}{根}。设 $n$ 是使 $(s_\alpha s_\beta)^n(\alpha)=\alpha$ 且 $(s_\alpha s_\beta)^n(\beta)=\beta$ 的最小正整数。则 $W$ 中由 $s_\alpha$ 与 $s_\beta$ 生成的\Term{subgroup}{子群} $W_{\alpha,\beta}$ 由关系
\[s_\alpha^2=1,\qquad s_\beta^2=1,\qquad(s_\alpha s_\beta)^n=1\]
定义。
\end{SGAresult}
考虑到 3.4.2，只须验证：
\begin{SGAresult}{引理}{3.4.4}
设 $G$ 是由两个元素 $x,y$ 按关系 $x^2=y^2=1$ 生成的群。$G$ 的任一既不包含 $x$ 也不包含 $y$ 的\Term{invariant-subgroup}{不变子群}，都由一个形如 $(xy)^n$（$n>0$）的元素生成（作为\Term{invariant-subgroup}{不变子群}）。
\end{SGAresult}
\begin{proof}
\footnote{编者注：下文对原文作了订正。}事实上，$G$ 的每个元素都可写成 $(xy)^n$，或 $(yx)^n=(xy)^{-n}$，或
\begin{align*}
\text{(a)}&\qquad x(yx)^{2n}\quad\text{或}\quad y(xy)^{2n+1},\\
\text{(b)}&\qquad y(xy)^{2n}\quad\text{或}\quad x(yx)^{2n+1},
\end{align*}
其中 $n\in\mathbb N$。而 (a) 型（分别地，(b) 型）元素与 $x$（分别地，$y$）\Term{conjugacy}{共轭}。
\end{proof}

\SourcePage{79}
\begin{SGAresult}{注记}{3.4.5}
整数 $n$ 可立即算出：若置
\[(\alpha^*,\beta)=p,\qquad(\beta^*,\alpha)=q,\]
则有\footnote{编者注：下文对原文作了订正。}
\[(s_\alpha s_\beta)(\alpha)=(pq-1)\alpha-q\beta,\qquad(s_\alpha s_\beta)(\beta)=p\alpha-\beta.\]
所以整数 $n$ 是矩阵
\[\begin{pmatrix}pq-1&p\\-q&-1\end{pmatrix}\]
的阶。

若 $pq=0$，则由 2.2.2，$p=q=0$，故 $n=2$。否则，由 2.3.1，$pq$ 等于 $1,2,3$ 中的一个，分别得到 $n=3,4,6$。

\textbf{注意。}写出上述矩阵的阶为\Term{finite}{有限}这一条件，可重新得到 2.3.1 中的不等式 (13)。
\end{SGAresult}
\begin{SGAresult}{定义}{3.4.6}
设 $\Delta$ 是一个\Term{simple-root-system}{单根系}，$\Delta'\subset\Delta$。记
\[R_{\Delta'}=R\cap(\mathbb Q\cdot S')=R\cap(\mathbb Z\cdot\Delta').\]
\end{SGAresult}
\begin{SGAresult}{引理}{3.4.7}
$R_{\Delta'}$ 是\Term{closed-root-set}{闭}且\Term{symmetric-root-set}{对称}的，$\Delta'$ 是\Term{root-datum}{根数据} $(M,M^*,R_{\Delta'},R_{\Delta'}^*)$ 的一个\Term{simple-root-system}{单根系}；其\Term{weyl-group}{Weyl 群}是 $W$ 中由 $s_\alpha$（$\alpha\in\Delta'$）生成的\Term{subgroup}{子群} $W_{\Delta'}$。有 $\Delta\cap R_{\Delta'}=\Delta'$。
\end{SGAresult}
显然。
\begin{SGAresult}{命题}{3.4.8}
设 $R'\subset R$。下列条件\Term{equivalence}{等价}：

(i) 存在 $V$（或 $\mathscr V(R)$）的一个向量子\Term{space}{空间} $V'$，使 $R'=R\cap V'$。

(ii) 存在 $R$ 的一个\Term{simple-root-system}{单根系} $\Delta$ 及 $\Delta$ 的一个\Term{subset}{子集} $\Delta'$，使 $R'=R_{\Delta'}$。

更准确地说，在这些条件下，$(M,M^*,R',R'^*)$ 的任一\Term{simple-root-system}{单根系} $\Delta'$ 都包含于 $R$ 的一个\Term{simple-root-system}{单根系} $\Delta$，并且 $R'=R_{\Delta'}$。
\end{SGAresult}
\begin{proof}
显然 (ii) $\Rightarrow$ (i)。假设 (i) 成立：则 $R'$ 是\Term{closed-root-set}{闭}且\Term{symmetric-root-set}{对称}的，故 $(M,M^*,R',R'^*)$ 是一个\Term{root-datum}{根数据}。设 $\Delta'$ 是此\Term{root-datum}{根数据}的一个\Term{simple-root-system}{单根系}，$R'_+=\mathscr P(\Delta')$。若 $V'=V$，则 $\Delta'$ 是 $R$ 的一个\Term{simple-root-system}{单根系}，结论已得。否则，存在 $x\in V^*$，使
\[(x,V')=\{0\},\qquad(x,\alpha)\ne0\quad\text{对每个 }\alpha\in R-R'.\]
置 $R_+(x)=\{\alpha\in R\mid(x,\alpha)>0\}$，$R_+=R_+(x)\cup R'_+$。对每个 $\alpha\in R$，有\Term{equivalence-relation}{等价关系}
\begin{align*}
(x,\alpha)>0&\Longleftrightarrow\alpha\in R_+(x),\\
(x,\alpha)<0&\Longleftrightarrow\alpha\in-R_+(x),\\
(x,\alpha)=0&\Longleftrightarrow\alpha\in R'.
\end{align*}
由 3.3.8 立即得到 $R_+$ 是 $R$ 的一个\Term{positive-root-system}{正根系}。置 $\Delta=\mathscr S(R_+)$。显然只须证明 $\Delta'\subset\Delta$。否则，取 $\alpha\in\Delta'-\Delta$。
\SourcePage{80}
由 3.2.14，存在 $\beta,\gamma\in R_+$ 使 $\alpha=\beta+\gamma$。若 $\beta,\gamma\in R_+(x)$，则 $\alpha\in R_+(x)$，这不可能，因为 $(x,S')=0$。若 $\beta$ 或 $\gamma$（例如 $\beta$）属于 $R'_+$，则由于 $R'$ 是\Term{symmetric-root-set}{对称}且\Term{closed-root-set}{闭}的，$\gamma=\alpha-\beta$ 属于 $R_+\cap R'=R'_+$；但这样 $\alpha$ 就不是 $R'_+$ 中的\Term{simple-root}{单根}。
\end{proof}
\begin{SGAresult}{引理}{3.4.9}
在上述条件下，设 $\alpha\in\mathscr P(\Delta)-R'$。对每个 $w\in W_{\Delta'}$，有 $w(\alpha)\in\mathscr P(\Delta)-R'$。
\end{SGAresult}
\begin{proof}
事实上，只须对 $w=s_\beta$（$\beta\in\Delta'$）验证；此时由 3.3.1 及 $s_\beta(R')=R'$ 即得。
\end{proof}
\begin{SGAresult}{引理}{3.4.10}
设 $w\in W$。在 3.4.8 的条件下，下列条件\Term{equivalence}{等价}：

(i) $w\in W_{\Delta'}$。

(ii) 对每个 $m\in M$，$w(m)-m\in V'$。

(iii) 对每个 $\alpha\in R$，$w(\alpha)-\alpha\in V'$。
\end{SGAresult}
\begin{proof}
显然有 (i) $\Rightarrow$ (ii) $\Rightarrow$ (iii)。证明 (iii) $\Rightarrow$ (i)。设 $w\in W$，且对每个 $\alpha\in R$ 有 $w(\alpha)-\alpha\in V'$。写成 $w=s_{\alpha_n}\cdots s_{\alpha_1}$，其中 $\alpha_i\in\Delta$，并对 $n$ 归纳证明每个 $\alpha_i\in\Delta'$。可假设 $w'=s_{\alpha_{n-1}}\cdots s_{\alpha_1}\in W_{\Delta'}$。于是对每个 $\alpha\in R$，
\[w(\alpha)-\alpha=w'(\alpha)-\alpha-(\alpha_n^*,w'(\alpha))\alpha_n.\]
取 $\alpha=w'^{-1}(\alpha_n)$，得到 $2\alpha_n\in V'$，故 $\alpha_n\in\Delta'$，所以 $w=s_{\alpha_n}w'$ 属于 $W_{\Delta'}$。
\end{proof}

\SGAsubsection{3.5}{其他注记}
\begin{SGAresult}{命题}{3.5.1}
设 $R_+$ 是一个\Term{positive-root-system}{正根系}。记
\[\rho_{R_+}=\frac12\sum_{\alpha\in\operatorname{ind}(R_+)}\alpha.\]
则对每个 $\beta\in\mathscr S(R_+)$ 都有 $(\beta^*,\rho_{R_+})=1$。
\end{SGAresult}
\begin{proof}
事实上，可写成
\[2\rho_{R_+}=\beta+\sum_{\substack{\alpha\in\operatorname{ind}(R_+)\\\alpha\ne\beta}}\alpha,\]
故由 3.3.2，$s_\beta(2\rho_{R_+})=2\rho_{R_+}-2\beta$。
\end{proof}
\begin{SGAresult}{推论}{3.5.2}
置 $\rho^*_{R_+}=\frac12\sum_{\alpha^*\in\operatorname{ind}(R_+^*)}\alpha^*$。则：

(i) 对每个 $\alpha\in R_+$，有 $(\rho^*_{R_+},\alpha)>0$（即 $\rho^*_{R_+}\in\mathscr C(R_+)$，见 3.6.8）。

(ii) 对每个 $\alpha\in R$，有 $\operatorname{ord}_{\mathscr S(R_+)}(\alpha)=(\rho^*_{R_+},\alpha)$。\footnote{编者注：事实上，由 3.1.5，只须对 $\alpha\in\Delta$ 验证此式。而把 3.5.1 应用于 $R^*$，此时得到 $(\rho^*_{R_+},\alpha)=1=\operatorname{ord}_{\mathscr S(R_+)}(\alpha)$。}
\end{SGAresult}

\SourcePage{81}
\begin{SGAresult}{注记}{3.5.3}
若 $w\in W$，则 $\rho_{w(R_+)}=w(\rho_{R_+})$，且 $\rho^*_{w(R_+)}=w(\rho^*_{R_+})$。
\end{SGAresult}
\begin{SGAresult}{命题}{3.5.4}
设 $\alpha$ 与 $\gamma$ 是两个不成比例的\Term{root}{根}，并假设 $\alpha$ \Term{indivisible-root}{不可分}。存在一个\Term{simple-root-system}{单根系}，其中包含 $\alpha$ 与一个\Term{root}{根} $\beta$，使 $\gamma=a\alpha+b\beta$，其中 $a,b\in\mathbb N$。
\end{SGAresult}
\begin{proof}
事实上，构造\Term{vector-space}{向量空间} $\mathscr V(R)$ 的一组基，其中包含 $\alpha_1=\alpha$、$\alpha_2=\gamma$。考虑关于这组基的\Term{lexicographic-order}{字典序}。把 $>0$ 的\Term{root}{根}的\Term{set}{集合}记为 $R_+$，显然 $R_+$ 是一个\Term{positive-root-system}{正根系}，且 $R_+$ 中不与 $\alpha$ 成比例的最小元素是\Term{simple-root}{单根}。\footnote{编者注：修改了原文，以涵盖 $2\alpha$ 是\Term{root}{根}的情形。}这个元素形如
\[\beta=p\alpha+q\gamma,\qquad0<q\leq1.\]
因此 $\gamma=q^{-1}\beta-q^{-1}p\alpha$，且因 $q^{-1}>0$，有 $q^{-1}\in\mathbb N^*$ 及 $-q^{-1}p\in\mathbb N$。
\end{proof}
最后对群 $\Gamma_0(R)$ 作两点注记。
\begin{SGAresult}{命题}{3.5.5}
设 $G$ 是一个\Term{abelian-group}{阿贝尔群}，$f:R\to G$ 是一个满足下列两个条件的映射：

(i) 若 $\alpha\in R$，则 $f(-\alpha)=-f(\alpha)$。

(ii) 若 $\alpha,\beta,\alpha+\beta\in R$，则 $f(\alpha+\beta)=f(\alpha)+f(\beta)$。

则存在唯一的群\Term{homomorphism}{同态} $\overline f:\Gamma_0(R)\to G$，使对 $\alpha\in R$ 有 $\overline f(\alpha)=f(\alpha)$。
\end{SGAresult}
\begin{proof}
事实上，若 $\Delta$ 是 $R$ 的一个\Term{simple-root-system}{单根系}，且 $\beta\in R$ 写成 $\sum_{\alpha\in\Delta}a(\alpha)\alpha$，则由 3.2.16 立即得 $f(\beta)=\sum_{\alpha\in\Delta}a(\alpha)f(\alpha)$。而 $\Delta$ 是 $\Gamma_0(R)$ 的一组基。
\end{proof}
\begin{SGAresult}{命题}{3.5.6}
设 $\Delta$ 是一个\Term{simple-root-system}{单根系}。$\Gamma_0(R)$ 上存在一个全序\Term{group-structure}{群结构}，使 $>0$ 的\Term{root}{根}正好是 $\mathscr P(\Delta)$ 的元素，且 $\alpha\mapsto\operatorname{ord}_\Delta(\alpha)$ 是递增函数。
\end{SGAresult}
\begin{proof}
事实上，设 $\alpha_1,\alpha_2,\ldots,\alpha_n$（$n=\operatorname{rgss}(\mathscr R)$）为 $\Delta$ 的各元素。对 $x\in\Gamma_0(R)$ 有\Term{resolution}{分解}
\[x=\sum_{i=1}^nm_i(x)\alpha_i.\]
只须取关于函数 $\sum m_i,m_n,m_{n-1},\ldots,m_2$ 的\Term{lexicographic-order}{字典序}。
\end{proof}
\begin{SGAresult}{注记}{3.5.7}
最前面的\Term{root}{根}依次为
\[\alpha_1,\alpha_2,\ldots,\alpha_n;\]
接下来是（若它们是\Term{root}{根}的话）$2\alpha_1,\alpha_1+\alpha_2,\alpha_1+\alpha_3,\ldots$。
\end{SGAresult}

\SourcePage{82}
\SGAsubsection{3.6}{Weyl 室}
\begin{SGAresult}{引理}{3.6.1}
设 $V$ 是一个\Term{finite}{有限}维 $\mathbb R$-\Term{vector-space}{向量空间}。\footnote{编者注：将 $\mathbb Q$ 换成了 $\mathbb R$。}设 $f_i$ 为\Term{linearly-independent}{线性无关}的\Term{linear-form}{线性形式}。置
\[C=\{x\in V\mid f_i(x)>0\}.\]
则 $C$ 是 $X=V-\bigcup_i f_i^{-1}(0)$ 的一个极大凸\Term{subset}{子集}。
\end{SGAresult}
显然。
\begin{SGAresult}{定义}{3.6.2}
$V$ 的一个可由 3.6.1 中方法描述的\Term{subset}{子集} $C$，（在这里）称为 $V$ 的一个\Term{chamber}{室}。
\end{SGAresult}
\begin{SGAresult}{定义}{3.6.3}
若 $V$ 的\Term{hyperplane}{超平面} $H$ 满足 $H\cap(\overline C-C)$ 包含 $H$ 的一个\Term{nonempty}{非空}\Term{open-subset}{开子集}，称 $H$ 为 $C$ 的一个\Term{chamber-wall}{壁}。
\end{SGAresult}
\begin{SGAresult}{注记}{3.6.4}
对于一个凸\Term{subset}{子集}，它的\Term{closure}{闭包}可不借助 $V$ 的\Term{topology}{拓扑}来描述：它是所有包含在给定\Term{subset}{子集}中的开线段的端点所成的\Term{set}{集合}。
\end{SGAresult}
\begin{SGAresult}{引理}{3.6.5}
在 3.6.1 的条件下，有
\[\overline C=\{x\in V\mid f_i(x)\geq0\}.\]
$C$ 的\Term{chamber-wall}{壁}是\Term{hyperplane}{超平面} $f_i^{-1}(0)$。
\end{SGAresult}
\begin{proof}
第一个断言显然。第二个由 $\overline C-C\subset\bigcup_i f_i^{-1}(0)$，以及各 $f_i^{-1}(0)$ 显然是 $C$ 的\Term{chamber-wall}{壁}得出。
\end{proof}
\begin{SGAresult}{命题}{3.6.6}
设 $C$ 是 $V$ 的一个\Term{chamber}{室}。若 $H_i$（$i=1,2,\ldots,n$）为 $C$ 的各个不同的\Term{chamber-wall}{壁}，则对任意满足 $H_i=u_i^{-1}(0)$ 的\Term{linear-form}{线性形式}组 $\{u_i\}$，存在 $\varepsilon_i\in\{-1,+1\}$，使 $C$ 由下式定义：
\[C=\{x\in V\mid\varepsilon_i u_i(x)>0\}.\]
对 $C$ 的任一个\Term{chamber-wall}{壁} $H$，有 $H\cap C=\varnothing$，且 $H\cap\overline C=\overline{C_H}$，其中 $C_H$ 是 $H$ 中的一个\Term{chamber}{室}。$C_{H_i}$ 的\Term{chamber-wall}{壁}为 $H_j\cap H_i$（$j\ne i$）。
\end{SGAresult}
这由引理显然得出。
\begin{SGAresult}{定义}{3.6.7}
各 $C_{H_i}$ 称为 $C$ 的\Term{chamber-face}{面}。
\end{SGAresult}
现在设 $\mathscr R=(M,M^*,R,R^*)$ 是一个\Term{root-datum}{根数据}。置 $V_{\mathbb R}^*=M^*\otimes\mathbb R$。
\begin{SGAresult}{定义}{3.6.8}
\footnote{编者注：补入了\Term{hyperplane}{超平面} $\mathscr H_\alpha$ 的定义。}对每个 $\alpha\in R$，置
\[\mathscr H_\alpha=\{x\in V_{\mathbb R}^*\mid(x,\alpha)=0\}.\]
记 $X=V_{\mathbb R}^*-\bigcup_{\alpha\in R}\mathscr H_\alpha$。对每个 $x\in X$，置
\[R_+(x)=\{\alpha\in R\mid(x,\alpha)>0\}.\]
\SourcePage{83}
对每个\Term{positive-root-system}{正根系} $R_+$，记
\[\mathscr C(R_+)=\{x\in V_{\mathbb R}^*\mid(x,\alpha)>0\text{ 对每个 }\alpha\in R_+\}.\]
\end{SGAresult}
\begin{SGAresult}{命题}{3.6.9}
(i) 对每个 $x\in X$，$R_+(x)$ 是一个\Term{positive-root-system}{正根系}。对每个\Term{positive-root-system}{正根系} $R_+$，$\mathscr C(R_+)$ 是 $V_{\mathbb R}^*$ 中的一个\Term{chamber}{室}。各 $\mathscr C(R_+)$ 是 $X$ 的极大凸\Term{subset}{子集}。

(ii) 设 $\Delta$ 是一个\Term{simple-root-system}{单根系}。有
\[\mathscr C(\mathscr P(\Delta))=\{x\in V_{\mathbb R}^*\mid(x,\alpha)>0\text{ 对每个 }\alpha\in\Delta\}.\]
$\mathscr C(\mathscr P(\Delta))$ 的\Term{chamber-wall}{壁}是\Term{hyperplane}{超平面} $\mathscr H_\alpha=\alpha^{-1}(0)$（$\alpha\in\Delta$）；它的\Term{chamber-face}{面}是
\[C_\alpha=\{x\in V_{\mathbb R}^*\mid(x,\alpha)=0,\ (x,\beta)>0\text{ 对 }\beta\in\Delta,\ \beta\ne\alpha\}.\]
(iii) 有\Term{equivalence-relation}{等价关系}
\[R_+(x)=R_+\quad\Longleftrightarrow\quad x\in\mathscr C(R_+).\]
\end{SGAresult}
\begin{proof}
首先，$R_+(x)$ 显然是一个\Term{positive-root-system}{正根系}。因 $x\in\mathscr C(R_+(x))$，各 $\mathscr C(R_+(x))$ 的并为 $X$。性质 (iii) 立即可见；由此得到各 $\mathscr C(R_+)$ 构成 $X$ 的一个分划。(ii) 的第一个断言显然。立即得出 $\mathscr C(R_+)$ 是 $V$ 中的一个\Term{chamber}{室}，这证明了 (ii) 的其余部分。最后只须注意
\[\bigcup_{\alpha\in R}\alpha^{-1}(0)=\bigcup_{\alpha\in\Delta}\alpha^{-1}(0),\]
便可由 3.6.1 完成 (i) 的证明。
\end{proof}
\begin{SGAresult}{定义}{3.6.10}
各 $\mathscr C(P)$ 称为此\Term{root-datum}{根数据}的\Term{weyl-chamber}{Weyl 室}。
\end{SGAresult}
\begin{SGAresult}{推论}{3.6.11}
映射 $R_+\mapsto\mathscr C(R_+)$ 与 $C\mapsto R_+(x)$（其中任取 $x\in C$），给出\Term{positive-root-system}{正根系}与\Term{weyl-chamber}{Weyl 室}之间的一个\Term{bijection}{双射}对应。
\end{SGAresult}
这个对应在\Term{weyl-group}{Weyl 群}\Term{action}{作用}下不变：
\begin{SGAresult}{引理}{3.6.12}
若 $w\in W(\mathscr R)$，则 $\mathscr C(w(R_+))=w(\mathscr C(R_+))$。
\end{SGAresult}
\begin{SGAresult}{推论}{3.6.13}
对应 $\Delta\leftrightarrow R_+\leftrightarrow C$ 是 $W(\mathscr R)$ 下的\Term{homogeneous-space}{齐性空间}的\Term{isomorphism}{同构}。
\end{SGAresult}
后文将看到，这些\Term{homogeneous-space}{齐性空间}是\Term{principal-homogeneous-space}{主齐性空间}（5.5）。
\begin{SGAresult}{注记}{3.6.14}
若 $C$ 是一个\Term{weyl-chamber}{Weyl 室}，则 $-C$ 也是一个\Term{weyl-chamber}{Weyl 室}，称为 $C$ 的\Term{opposite-chamber}{对室}。因此存在 $w_0\in W$（事实上只有一个，见 5.5），使 $w_0(C)=-C$；称其为此\Term{root-datum}{根数据}关于\Term{weyl-chamber}{Weyl 室} $C$（或关于 $\mathscr P(C)$，或关于 $\mathscr S(\mathscr P(C))$，等等）的\Term{root-datum-symmetry}{对称变换}。
\end{SGAresult}

\SourcePage{84}
\SGAsection{4}{半单秩为 2 的约化根数据}
\textbf{4.0。}\footnote{编者注：增加了编号 4.0，以供后文引用。}设 $\mathscr R$ 是\Term{semisimple-rank}{半单秩}为 $2$ 的一个\Term{root-datum}{根数据}。设 $\{\alpha,\beta\}$ 是一个\Term{simple-root-system}{单根系}。假设 $\ell(\alpha)\leq\ell(\beta)$。由 2.3.1 与 3.2.1，有下列四种可能：
\[
\begin{array}{c|cccc}
\text{类型}&\ell(\beta)/\ell(\alpha)&\ell(\beta^*)/\ell(\alpha^*)&(\beta^*,\alpha)&(\alpha^*,\beta)\\\hline
A_1\times A_1&-&-&0&0\\
A_2&1&1&-1&-1\\
B_2&2&1/2&-1&-2\\
G_2&3&1/3&-1&-3
\end{array}
\]
于是由 3.4.5，$s_\alpha s_\beta$ 的阶分别为 $2,3,4,6$。

分别研究这些\Term{root-system}{根系}，并列出\Term{indivisible-root}{不可分}\Term{root}{根}。

\textbf{$A_1\times A_1$ 型。}\Term{indivisible-root}{不可分}\Term{root}{根}为 $\alpha,\beta,-\alpha,-\beta$。对应的\Term{coroot}{余根}为 $\alpha^*,\beta^*,-\alpha^*,-\beta^*$。

\textbf{$A_2$ 型。}\Term{indivisible-root}{不可分}\Term{positive-root}{正根}如下：
\[
\begin{array}{c|ccc}
\text{根 }\gamma&\alpha&\beta&\alpha+\beta\\\hline
\ell(\gamma)/\ell(\alpha)&1&1&1\\\hline
\text{余根 }\gamma^*&\alpha^*&\beta^*&\alpha^*+\beta^*\\\hline
\ell(\gamma^*)/\ell(\beta^*)&1&1&1
\end{array}
\]
\Term{indivisible-root}{不可分}\Term{positive-root}{正根}的半和是 $\rho=\alpha+\beta$。

\textbf{$B_2$ 型。}\Term{indivisible-root}{不可分}\Term{positive-root}{正根}如下：
\[
\begin{array}{c|cccc}
\text{根 }\gamma&\alpha&\beta&\alpha+\beta&2\alpha+\beta\\\hline
\ell(\gamma)/\ell(\alpha)&1&2&1&2\\\hline
\text{余根 }\gamma^*&\alpha^*&\beta^*&2\alpha^*+\beta^*&\alpha^*+\beta^*\\\hline
\ell(\gamma^*)/\ell(\beta^*)&2&1&2&1
\end{array}
\]
\Term{indivisible-root}{不可分}\Term{positive-root}{正根}的半和是 $\rho=(4\alpha+3\beta)/2$。

\textbf{$G_2$ 型。}\Term{indivisible-root}{不可分}\Term{positive-root}{正根}如下：
\[
\begin{array}{c|cccccc}
\text{根 }\gamma&\alpha&\beta&\alpha+\beta&2\alpha+\beta&3\alpha+\beta&3\alpha+2\beta\\\hline
\ell(\gamma)/\ell(\alpha)&1&3&1&1&3&3\\\hline
\text{余根 }\gamma^*&\alpha^*&\beta^*&\alpha^*+3\beta^*&2\alpha^*+3\beta^*&\alpha^*+\beta^*&\alpha^*+2\beta^*\\\hline
\ell(\gamma^*)/\ell(\beta^*)&3&1&3&3&1&1
\end{array}
\]
\Term{indivisible-root}{不可分}\Term{positive-root}{正根}的半和是 $\rho=5\alpha+3\beta$。

\SourcePage{85}
\begin{SGAresult}{命题}{4.1}
\begingroup\renewcommand{\thefootnote}{*}\footnote[0]{下列 4.1 至 4.4 用于 5.1 的证明。现在已有更简单的 5.1 的证明（见 [BLie]，第 V.3 节，定理 1）。编者注：明确了参考文献，并把此注移至此处；原文中它位于 4.2。}\endgroup 设 $n$ 是 $s_\alpha s_\beta$ 的阶。置 $u_0=0$，并对 $p\in\mathbb N$，置
\[
\begin{cases}
u_{2p+1}=u_{2p}+(s_\alpha s_\beta)^p(\alpha);\\
u_{2p+2}=u_{2p+1}+(s_\alpha s_\beta)^ps_\alpha(\beta).
\end{cases}
\]
则：\footnote{编者注：设 $\rho$ 是\Term{positive-root}{正根}的半和（见 3.5.1）；若像下文 4.2 那样置 $w_0=1,w_1=s_\alpha,w_2=s_\beta s_\alpha,w_3=s_\alpha s_\beta s_\alpha$，等等，则 $u_k$ 就是 $\rho-w_k^{-1}(\rho)$，由 $w_{2n}=\operatorname{id}$ 即得 (i)。}

(i) 对每个 $k\in\mathbb N$，$u_{k+2n}=u_k$。

(ii) $u_0=0$，$u_1=\alpha$，$u_{2n-1}=\beta$，$u_{2n}=0$。

(iii) 若 $1<k<2n-1$，则 $u_k=a_k\alpha+b_k\beta$，其中 $a_k,b_k\in\mathbb N^*$。
\end{SGAresult}
\begin{proof}
断言 (i) 由 $(s_\alpha s_\beta)^n=1$ 与 $u_{2n}=0$ 得出。

证明 (ii) 与 (iii)。在这四种情形中，计算立即给出下列数值序列：\footnote{编者注：订正了原文，使之符合 4.0 中 $\ell(\alpha)\leq\ell(\beta)$ 的约定。}
\begin{align*}
(A_1\times A_1)\quad&0,\alpha,\beta+\alpha,\beta,0.\\[2pt]
(A_2)\quad&0,\alpha,\beta+2\alpha,2\beta+2\alpha,2\beta+\alpha,\beta,0.\\[2pt]
(B_2)\quad&0,\alpha,\beta+3\alpha,2\beta+4\alpha,3\beta+4\alpha,\\
&3\beta+3\alpha,2\beta+\alpha,\beta,0.\\[2pt]
(G_2)\quad&0,\alpha,\beta+4\alpha,2\beta+6\alpha,4\beta+9\alpha,5\beta+10\alpha,\\
&6\beta+10\alpha,6\beta+9\alpha,5\beta+6\alpha,4\beta+4\alpha,2\beta+\alpha,\beta,0.
\end{align*}
\end{proof}
\begin{SGAresult}{引理}{4.2}
置 $w_{2p}=(s_\beta s_\alpha)^p$，$w_{2p+1}=s_\alpha(s_\beta s_\alpha)^p$，这样 $w_0,\ldots,w_{2n-1}$ 就是 $W$ 的各个不同的元素。设 $u_0,\ldots,u_{2n-1}$ 是 4.1 中定义的 $V$ 的元素。对每个 $x\in V^*$，有
\[x-w_k(x)=n_k\alpha^*+m_k\beta^*,\]
其中 $n_k,m_k\in\mathbb Q$，且 $n_k+m_k=(x,u_k)$。\footnote{编者注：考虑到编者注 (17)，这由 3.5.1 与 1.1.9 立即得出：$n_k+m_k=(x-w_k(x),\rho)=(x,\rho-w_k^{-1}(\rho))=(x,u_k)$。}
\end{SGAresult}
\begin{proof}
对 $k$ 归纳。若 $k=0$，公式显然成立。例如，由 $w_{2p}$ 过渡到 $w_{2p+1}$。因 $w_{2p+1}=s_\alpha w_{2p}$，故
\begin{align*}
x-w_{2p+1}(x)&=x-w_{2p}(x)+w_{2p}(x)-s_\alpha w_{2p}(x)\\
&=n_{2p}\alpha^*+m_{2p}\beta^*+(w_{2p}(x),\alpha)\alpha^*.
\end{align*}
所以
\begin{align*}
n_{2p+1}+m_{2p+1}&=n_{2p}+m_{2p}+((s_\beta s_\alpha)^p(x),\alpha)\\
&=(x,u_{2p})+(x,(s_\alpha s_\beta)^p(\alpha))=(x,u_{2p+1}).
\end{align*}
\end{proof}

\SourcePage{86}
\begin{SGAresult}{推论}{4.3}
\footnote{编者注：订正了命题陈述。}设 $x\in V^*$。对每个 $w\in W$，置
\[x-w(x)=a_w\alpha^*+b_w\beta^*.\]
若 $(x,\alpha)\geq0$ 且 $(x,\beta)\geq0$，则 $a_w+b_w\geq0$。若还有 $(x,\alpha)>0$（分别地，$(x,\beta)>0$），则当 $w\ne1,s_\beta$（分别地，当 $w\ne1,s_\alpha$）时，$a_w+b_w>0$。
\end{SGAresult}
这由 4.1 与 4.2 立即得出。
\begin{SGAresult}{推论}{4.4}
设 $\mathscr R$ 是任意一个\Term{root-datum}{根数据}，$\Delta$ 是一个\Term{simple-root-system}{单根系}。设 $\gamma$ 是一个\Term{positive-root}{正根}，$\alpha,\beta$ 是两个\Term{simple-root}{单根}，$W_{\alpha,\beta}$ 是 $W$ 中由 $s_\alpha$ 与 $s_\beta$ 生成的\Term{subgroup}{子群}。若
\[\operatorname{ord}_\Delta(s_\alpha(\gamma))<\operatorname{ord}_\Delta(\gamma),\qquad
\operatorname{ord}_\Delta(s_\beta(\gamma))\leq\operatorname{ord}_\Delta(\gamma),\]
则对每个 $w\in W_{\alpha,\beta}$，有 $\operatorname{ord}_\Delta(w(\gamma))\leq\operatorname{ord}_\Delta(\gamma)$；此外，若 $w\ne1$、$w\ne s_\beta$，则 $\operatorname{ord}_\Delta(w(\gamma))<\operatorname{ord}_\Delta(\gamma)$。
\end{SGAresult}
\begin{proof}
事实上，考虑\Term{dual-root-datum}{对偶根数据} $\mathscr R^*$，再考虑\Term{root-datum}{根数据} $(M^*,M,R'^*,R')$，其中 $R'$ 为所有能写成 $\alpha$ 与 $\beta$ 的有理系数线性组合的\Term{root}{根}的\Term{set}{集合}。对这个\Term{root-datum}{根数据}应用 4.3，即得所述推论。\footnote{编者注：事实上，假设\Term{equivalence}{等价}于 $(\gamma,\alpha^*)>0$ 且 $(\gamma,\beta^*)\geq0$；于是，对每个 $w\in W_{\alpha,\beta}$，$\gamma-w(\gamma)$ 属于 $\mathbb N\alpha+\mathbb N\beta$，且当 $w\notin\{1,s_\beta\}$ 时它不为零。}
\end{proof}
